\section{VERIFICAÇÃO, VALIDAÇÃO E TESTE}

\subsection{Parte 1: Estratégia de Testes}

A estratégia de testes do MedAgenda é dividida em quatro níveis fundamentais para garantir o correto funcionamento e a qualidade do software \cite{myers2011}:

\subsubsection{1. Testes de Unidade}
\begin{itemize}
    \item \textbf{Objetivo}: Validar a precisão de funções e métodos isolados da aplicação.
    \item \textbf{O que será testado}: Validação de formato de e-mail e CPF, e verificação da regra de antecedência mínima de 24 horas no cancelamento.
    \item \textbf{Justificativa}: Permite detectar falhas lógicas rapidamente e com baixo custo de correção.
\end{itemize}

\subsubsection{2. Testes de Integração}
\begin{itemize}
    \item \textbf{Objetivo}: Verificar a comunicação correta entre os serviços do sistema e o banco de dados.
    \item \textbf{O que será testado}: Integração entre \texttt{ConsultaService} e o banco de dados PostgreSQL, além da comunicação com o serviço de e-mail.
    \item \textbf{Justificativa}: Garante que os módulos colaboram sem erros de inconsistência de dados.
\end{itemize}

\subsubsection{3. Testes de Sistema}
\begin{itemize}
    \item \textbf{Objetivo}: Avaliar o funcionamento global da aplicação completa em ambiente homologado.
    \item \textbf{O que será testado}: Fluxo completo de agendamento, cancelamento, login e emissão de relatórios.
    \item \textbf{Justificativa}: Comprova que o sistema cumpre todos os requisitos funcionais e não funcionais contratados.
\end{itemize}

\subsubsection{4. Testes de Aceitação}
\begin{itemize}
    \item \textbf{Objetivo}: Validar se o sistema atende adequadamente às rotinas diárias dos usuários finais.
    \item \textbf{O que será testado}: Facilidade de agendamento por pacientes e clareza da agenda visual dos médicos.
    \item \textbf{Justificativa}: Assegura a utilidade prática e a satisfação do usuário.
\end{itemize}

\subsection{Parte 2: Casos de Teste}

A Tabela \ref{tab:casos_teste} apresenta os casos de teste elaborados para os requisitos funcionais prioritários do sistema.

\begin{table}[H]
\centering
\caption{Casos de Teste prioritários do sistema MedAgenda.}
\label{tab:casos_teste}
\begin{tabular}{|c|c|p{3.5cm}|p{3.5cm}|p{3.8cm}|}
\hline
\textbf{ID} & \textbf{RF} & \textbf{Objetivo} & \textbf{Entrada} & \textbf{Resultado Esperado} \\
\hline
CT-01 & RF-01 & Validar cadastro de paciente & Nome, CPF e e-mail válidos & Cadastro salvo com sucesso e confirmação exibida. \\
\hline
CT-02 & RF-02 & Validar bloqueio de login incorreto & E-mail correto e senha errada & Login negado e mensagem de erro exibida. \\
\hline
CT-03 & RF-03 & Validar agendamento de consulta & Médico, data futura e horário livre & Consulta agendada com status AGENDADA. \\
\hline
CT-04 & RF-04 & Validar regra de cancelamento & Tentativa de cancelar com menos de 24h & Cancelamento bloqueado com aviso da regra. \\
\hline
CT-05 & RF-05 & Validar visualização da agenda & Médico acessa agenda do dia & Exibição correta das consultas marcadas. \\
\hline
CT-06 & RF-06 & Validar registro de prontuário & Preenchimento de anamnese e diagnóstico & Prontuário salvo e consulta concluída. \\
\hline
CT-07 & RF-08 & Validar envio de lembrete & Simulação de 24h antes da consulta & E-mail de lembrete enviado ao paciente. \\
\hline
CT-08 & RF-11 & Validar histórico do paciente & Acesso à aba Meu Histórico & Listagem cronológica das consultas realizadas. \\
\hline
\end{tabular}\\
\autoria
\end{table}

\subsection{Parte 3: Matriz de Rastreabilidade de Requisitos para Casos de Teste}

A Tabela \ref{tab:matriz_testes} relaciona cada requisito funcional ao seu respectivo caso de teste.

\begin{table}[H]
\centering
\caption{Matriz de Rastreabilidade Requisitos $\times$ Casos de Teste.}
\label{tab:matriz_testes}
\begin{tabular}{|c|l|c|}
\hline
\textbf{Requisito Funcional (RF)} & \textbf{Descrição Resumida} & \textbf{Caso de Teste} \\
\hline
RF-01 & Cadastro de Paciente & CT-01 \\
\hline
RF-02 & Autenticação de Usuários & CT-02 \\
\hline
RF-03 & Agendamento de Consulta & CT-03 \\
\hline
RF-04 & Cancelamento de Consulta & CT-04 \\
\hline
RF-05 & Visualização de Agenda Médica & CT-05 \\
\hline
RF-06 & Registro de Prontuário & CT-06 \\
\hline
RF-08 & Notificação de Lembrete & CT-07 \\
\hline
RF-11 & Visualização de Histórico & CT-08 \\
\hline
\end{tabular}\\
\autoria
\end{table}
